\documentclass[11pt,a4paper]{article}

\usepackage[utf8]{inputenc}
\usepackage[T1]{fontenc}
\usepackage[spanish]{babel}
\usepackage[margin=2.5cm]{geometry}
\usepackage{hyperref}
\usepackage{listings}
\usepackage{xcolor}
\usepackage{booktabs}
\usepackage{enumitem}

\hypersetup{
    colorlinks=true,
    linkcolor=blue,
    urlcolor=blue,
    citecolor=blue
}

\lstset{
    basicstyle=\ttfamily\small,
    breaklines=true,
    frame=single,
    columns=fullflexible,
    backgroundcolor=\color{gray!10}
}

\title{Informe técnico\\
Sistema de gestión de inventario del Mercado Municipal de Quevedo}
\author{Loor Marlon \\ Universidad Técnica Estatal de Quevedo (UTEQ) \\ Aplicaciones Web}
\date{\today}

\begin{document}

\maketitle

\begin{abstract}
Este informe describe el diseño técnico del sistema de gestión de inventario
desarrollado para el examen parcial de Aplicaciones Web: una API REST construida
con Spring Boot 3.5.16 sobre Java 21, con autenticación y autorización basadas en
JSON Web Tokens (JWT), persistencia en PostgreSQL, y una capa de cache con Redis
bajo el patrón \emph{cache-aside}. Se documentan la arquitectura por capas, el
contrato de respuesta uniforme de la API, el modelo de datos, y las decisiones de
seguridad y de caché tomadas durante la implementación.
\end{abstract}

\section{Introducción}

El sistema permite gestionar el inventario de productos del Mercado Municipal de
Quevedo mediante una API REST protegida por autenticación JWT. Los requisitos
funcionales cubiertos son: listado paginado de productos, creación validada,
eliminación lógica (\emph{soft delete}), autorización basada en roles
(\texttt{ROLE\_USER} / \texttt{ROLE\_ADMIN}), y una capa de cache con Redis para
reducir la carga sobre la base de datos en las operaciones de lectura.

Todo el sistema se distribuye mediante Docker Compose, con tres servicios:
la aplicación Spring Boot (\texttt{app}), PostgreSQL 16 (\texttt{postgres}) y
Redis 7 (\texttt{redis}).

\section{Arquitectura por capas}

El proyecto sigue una arquitectura por capas convencional, separando
responsabilidades para que cada componente tenga un único motivo de cambio:

\begin{itemize}[noitemsep]
    \item \textbf{\texttt{domain}}: entidades JPA (\texttt{Producto},
        \texttt{Usuario}) que mapean directamente las tablas de la base de datos.
    \item \textbf{\texttt{repository}}: interfaces \texttt{JpaRepository}
        (\texttt{ProductoRepository}, \texttt{UsuarioRepository}). Es el
        \emph{único} punto de acceso a datos: ni los controladores ni los
        servicios ejecutan consultas directas.
    \item \textbf{\texttt{service}}: lógica de negocio y orquestación de caché
        (\texttt{ProductoService}). Traduce entidades a DTOs de salida.
    \item \textbf{\texttt{controller}}: adaptadores HTTP delgados
        (\texttt{ProductoController}, \texttt{AuthController}) que solo mapean
        peticiones/respuestas y delegan toda la lógica al servicio.
    \item \textbf{\texttt{dto}}: objetos de transferencia (\texttt{ProductoRequest},
        \texttt{ProductoResponse}, \texttt{ApiResponse}, etc.), implementados como
        \emph{records} de Java 21 por su inmutabilidad y menor código repetitivo.
    \item \textbf{\texttt{security}}: todo lo relacionado con JWT y Spring
        Security (\texttt{JwtService}, \texttt{JwtAuthenticationFilter},
        \texttt{SecurityConfig}, manejadores de errores de autenticación).
    \item \textbf{\texttt{web}}: manejo centralizado de excepciones
        (\texttt{GlobalExceptionHandler}) que traduce cualquier excepción de la
        capa de negocio al contrato de respuesta uniforme de la API.
    \item \textbf{\texttt{config}}: configuración de infraestructura transversal,
        como la caché de Redis (\texttt{CacheConfig}).
\end{itemize}

Esta separación seaproxima a los principios de diseño de APIs orientadas a
recursos descritos por Fielding~\cite{fielding2000rest}, en cuanto a que cada
recurso (\texttt{/api/v1/productos}) se manipula mediante los verbos HTTP
estándar (\texttt{GET}, \texttt{POST}, \texttt{DELETE}) sin exponer detalles de
implementación interna al cliente.

\section{Modelo de datos}

El esquema (\texttt{db/schema.sql}) define dos tablas:

\begin{lstlisting}[language=SQL]
CREATE TABLE productos (
    id BIGSERIAL PRIMARY KEY,
    nombre VARCHAR(100) NOT NULL,
    categoria VARCHAR(50) NOT NULL,
    stock INTEGER NOT NULL CHECK (stock >= 0),
    precio DECIMAL(10,2) NOT NULL CHECK (precio >= 0.01),
    activo BOOLEAN NOT NULL DEFAULT TRUE,
    creado_en TIMESTAMPTZ NOT NULL DEFAULT now()
);

CREATE TABLE usuarios (
    id BIGSERIAL PRIMARY KEY,
    username VARCHAR(50) NOT NULL UNIQUE,
    password_hash VARCHAR(100) NOT NULL,
    role VARCHAR(20) NOT NULL CHECK (role IN ('ROLE_USER', 'ROLE_ADMIN'))
);
\end{lstlisting}

Los \texttt{CHECK} de \texttt{stock} y \texttt{precio} garantizan invariantes de
negocio a nivel de base de datos (no solo en la capa de validación de la API),
y el borrado de productos es siempre lógico: el campo \texttt{activo} se pone en
\texttt{false} en lugar de ejecutar un \texttt{DELETE} físico, preservando el
historial. El listado solo considera productos con \texttt{activo = true}
(\texttt{ProductoRepository.findByActivoTrue}).

Las contraseñas de \texttt{usuarios} se almacenan con \texttt{BCrypt}
(\texttt{password\_hash}), siguiendo el esquema adaptativo de
Provos y Mazières~\cite{provos1999bcrypt}, que incorpora un factor de costo
configurable para resistir ataques de fuerza bruta a medida que el hardware
mejora.

\section{Contrato de respuesta uniforme}

Todas las respuestas de la API, exitosas o de error, siguen la misma forma,
implementada como el \emph{record} genérico \texttt{ApiResponse<T>}:

\begin{lstlisting}[language=json]
{
  "success": true,
  "data": [ ... ],
  "message": "Listado obtenido correctamente",
  "meta": { "page": 0, "size": 20, "totalElements": 57, "totalPages": 3 }
}
\end{lstlisting}

En una respuesta de error, el campo \texttt{errors} (lista de
\texttt{\{field, message\}}) reemplaza a \texttt{data}/\texttt{meta}:

\begin{lstlisting}[language=json]
{
  "success": false,
  "message": "Error de validacion",
  "errors": [
    { "field": "nombre", "message": "must not be blank" }
  ]
}
\end{lstlisting}

Se optó por una \'unica clase \texttt{ApiResponse<T>} (con \texttt{meta} y
\texttt{errors} anotados con \texttt{@JsonInclude(NON\_NULL)}) en lugar de una
jerarquía de tipos de éxito/error separada, de modo que todos los controladores
y el manejador global de excepciones devuelvan siempre el mismo tipo, y el
cliente de la API pueda inspeccionar únicamente el campo \texttt{success} para
decidir cómo interpretar la respuesta.

\section{Seguridad: autenticación y autorización con JWT}

La autenticación es \emph{stateless}: el cliente obtiene un token mediante
\texttt{POST /api/v1/auth/login} y lo envía en cada petición protegida con el
esquema \texttt{Authorization: Bearer <token>}, tal como lo define
RFC 6750~\cite{rfc6750} para el uso de tokens portador (\emph{bearer tokens}).

El token en sí es un JSON Web Token firmado con HMAC-SHA256, según
RFC 7519~\cite{rfc7519}, con dos \emph{claims} relevantes además de los
estándar (\texttt{sub}, \texttt{iat}, \texttt{exp}): el nombre de usuario como
\emph{subject} y el rol (\texttt{ROLE\_USER} o \texttt{ROLE\_ADMIN}) como claim
personalizado \texttt{role}. El tiempo de expiración se fijó en una hora,
un valor razonable para pruebas y demostraciones que limita la ventana de
exposición si el token se filtra, sin forzar reautenticaciones constantes.

El flujo de autorización sigue el modelo estándar de Spring
Security~\cite{springsecurity2024}:

\begin{enumerate}[noitemsep]
    \item \texttt{JwtAuthenticationFilter} (un \texttt{OncePerRequestFilter})
        intercepta cada petición, extrae el token del header
        \texttt{Authorization}, y si es válido puebla el
        \texttt{SecurityContext} con las \emph{authorities} correspondientes.
    \item Si no hay token o es inválido, el filtro simplemente deja continuar
        la cadena sin autenticar; es \texttt{AuthorizationFilter} quien más
        adelante rechaza el acceso a un recurso protegido, delegando el error
        \texttt{401} a un \texttt{AuthenticationEntryPoint} personalizado.
    \item Un \texttt{AccessDeniedHandler} personalizado traduce los rechazos
        por rol insuficiente (\texttt{403}) al mismo contrato
        \texttt{ApiResponse.error(...)} usado en el resto de la API, en vez del
        HTML por defecto de Spring Security.
\end{enumerate}

Las reglas de autorización, definidas en \texttt{SecurityConfig}, son:

\begin{center}
\begin{tabular}{@{}lll@{}}
\toprule
Endpoint & Rol requerido & Sin token / rol insuficiente \\
\midrule
\texttt{POST /api/v1/auth/login} & público & --- \\
\texttt{GET /api/v1/productos} & \texttt{ROLE\_USER} o \texttt{ROLE\_ADMIN} & 401 \\
\texttt{POST /api/v1/productos} & \texttt{ROLE\_ADMIN} & 401 / 403 \\
\texttt{DELETE /api/v1/productos/\{id\}} & \texttt{ROLE\_ADMIN} & 401 / 403 \\
\bottomrule
\end{tabular}
\end{center}

Dado que Spring Security no hace que \texttt{ROLE\_ADMIN} herede automáticamente
las autoridades de \texttt{ROLE\_USER}, el endpoint de listado se declaró
explícitamente con \texttt{hasAnyRole("USER", "ADMIN")} en lugar de introducir
una jerarquía de roles (\texttt{RoleHierarchy}) separada: con un único rol por
usuario en el modelo de datos, esta forma es más directa de auditar en un solo
lugar.

Como salvaguarda adicional, la construcción de \texttt{GrantedAuthority} a
partir del valor de rol (ya sea leído de la base de datos o de un claim del
JWT) se centralizó en un único método utilitario (\texttt{RoleAuthorityUtils})
que normaliza el prefijo \texttt{ROLE\_} si no está presente, evitando que una
futura fuente de datos sin ese prefijo rompa la autorización en silencio.

\section{Patrón cache-aside con Redis}

El listado paginado de productos (\texttt{GET /api/v1/productos}) usa el patrón
\emph{cache-aside}, tal como lo describen tanto la documentación de patrones de
Azure~\cite{msftcacheaside} como Kleppmann~\cite{kleppmann2017ddia} en su
discusión sobre sistemas de caché: la aplicación consulta primero la caché: si
el valor no está (\emph{cache miss}), calcula el resultado consultando la base
de datos, lo guarda en la caché, y lo devuelve; en peticiones subsecuentes con
los mismos parámetros de paginación (\emph{cache hit}), el resultado se sirve
directamente desde Redis sin volver a consultar PostgreSQL.

En la implementación, esto se logra anotando el método de servicio (no el
controlador) con \texttt{@Cacheable}:

\begin{lstlisting}[language=java]
@Cacheable(
    value = "productos",
    key = "'page:' + #pageable.pageNumber
         + ':size:' + #pageable.pageSize
         + ':sort:' + #pageable.sort.toString()"
)
public ProductoPageResult listarActivos(Pageable pageable) { ... }
\end{lstlisting}

La clave de caché incluye página, tamaño y orden para no mezclar resultados de
distintas combinaciones de paginación bajo la misma entrada. Cualquier
operación de escritura (\texttt{crear}, \texttt{eliminar}) invalida toda la
caché de productos con \texttt{@CacheEvict(value = "productos", allEntries =
true)}, dado que una única entrada modificada puede afectar el conteo total y
la distribución de varias páginas cacheadas simultáneamente.

Para la serialización se usó un \texttt{Jackson2JsonRedisSerializer} atado al
tipo concreto \texttt{ProductoPageResult} (en vez del serializador genérico con
tipado polimórfico de Spring Data Redis), ya que los DTOs del proyecto son
\emph{records} de Java --- por lo tanto, clases \texttt{final} --- para las que
la estrategia de tipado por defecto de Jackson omite la metadata de tipo en el
valor raíz, produciendo un documento que el propio Jackson no puede releer. Al
conocerse de antemano el único tipo que esta caché almacena, un serializador
tipado explícito es más simple y evita ese problema por completo.

\section{Manejo centralizado de errores}

\texttt{GlobalExceptionHandler} (\texttt{@RestControllerAdvice}) concentra la
traducción de excepciones al contrato \texttt{ApiResponse.error(...)}:

\begin{itemize}[noitemsep]
    \item \texttt{AuthenticationException} (credenciales inválidas en el
        login) $\rightarrow$ \texttt{401}.
    \item \texttt{MethodArgumentNotValidException} (fallos de validación de
        \texttt{@Valid}) $\rightarrow$ \texttt{400}, con la lista de errores
        por campo.
    \item \texttt{ProductoNoEncontradoException} (id inexistente en
        \texttt{DELETE}) $\rightarrow$ \texttt{404}.
\end{itemize}

Centralizar este mapeo en un único componente evita duplicar la construcción
de la respuesta de error en cada controlador, y mantiene el formato de salida
consistente en toda la API.

\section{Conclusiones}

El sistema implementado cubre los cinco requisitos funcionales del examen
(listado paginado, creación validada, eliminación lógica, autorización basada
en roles, y caché con Redis) sobre una arquitectura por capas convencional que
aísla claramente el acceso a datos, la lógica de negocio y la capa HTTP. Las
decisiones de seguridad (JWT stateless, BCrypt, normalización centralizada de
roles) y de caché (cache-aside con invalidación total en escritura, y
serialización tipada para evitar los problemas de tipado polimórfico de
Jackson con \emph{records}) se documentaron junto con su justificación técnica,
de forma que puedan auditarse y, si corresponde, revisarse en iteraciones
futuras del proyecto.

\bibliographystyle{plain}
\bibliography{referencias}

\end{document}
